\begin{tabularx}{\textwidth}{L{25mm}L{32mm}X}
\toprule
\textbf{Záradék} & \textbf{Relációs algebrai szerep} & \textbf{Feladat} \\
\midrule
\code{FROM} & Operandusrelációk kijelölése & Megadja, mely táblákból, nézetekből vagy részlekérdezésekből indul a lekérdezés. \\
\code{JOIN ... ON} & Join, Descartes-szorzat szűréssel & Összekapcsolja a forrásokat, gyakran elsődleges kulcs--idegen kulcs kapcsolat alapján. \\
\code{WHERE} & Szelekció & Rekordszintű feltétellel szűri az eredményt a csoportosítás előtt. \\
\code{GROUP BY} & Csoportképzés & A rekordokat csoportokra bontja, hogy csoportonként lehessen aggregálni. \\
\code{HAVING} & Csoportokra vonatkozó szelekció & Aggregált értékek alapján szűri a csoportokat. \\
\code{SELECT} & Projekció és kiterjesztés & Meghatározza, mely mezők és számított kifejezések kerüljenek az eredménybe. \\
\code{ORDER BY} & Rendezés, nem tisztán algebrai művelet & Megadja az eredmény megjelenítési sorrendjét; a relációs algebra halmazszemlélete miatt ez külön SQL-funkció. \\
\bottomrule
\end{tabularx}
